\documentclass[12pt]{article}

% --- PACKAGES ---
\usepackage[utf8]{inputenc}
\usepackage[a4paper, margin=1in]{geometry}
\usepackage{graphicx}
\usepackage{amsmath}
\usepackage{booktabs}
\usepackage{hyperref}
\usepackage{float}
\usepackage{enumitem}
\usepackage{array}
\usepackage{listings}
\usepackage{xcolor}
\usepackage{mwe} % For placeholder images, you can remove this

% --- HYPERLINK AND LISTING STYLES ---
\hypersetup{
    colorlinks=true,
    linkcolor=blue,
    filecolor=magenta,      
    urlcolor=cyan,
}

\definecolor{codegreen}{rgb}{0,0.6,0}
\definecolor{codegray}{rgb}{0.5,0.5,0.5}
\definecolor{codepurple}{rgb}{0.58,0,0.82}
\definecolor{backcolour}{rgb}{0.95,0.95,0.92}

\lstdefinestyle{mystyle}{
    backgroundcolor=\color{backcolour},   
    commentstyle=\color{codegreen},
    keywordstyle=\color{magenta},
    numberstyle=\tiny\color{codegray},
    stringstyle=\color{codepurple},
    basicstyle=\footnotesize\ttfamily,
    breakatwhitespace=false,         
    breaklines=true,                 
    captionpos=b,                    
    keepspaces=true,                 
    numbers=left,                    
    numbersep=5pt,                  
    showspaces=false,                
    showstringspaces=false,
    showtabs=false,                  
    tabsize=2
}
\lstset{style=mystyle}

%====================================================================
% DOCUMENT START
%====================================================================
\begin{document}

% --- TITLE PAGE ---
\begin{center}
    \large \textbf{Sri Sivasubramaniya Nadar College of Engineering, Chennai} \\
    (An autonomous Institution affiliated to Anna University) \\
    \vspace{1cm}
\end{center}

\begin{table}[H]
\renewcommand{\arraystretch}{1.5}
\resizebox{\textwidth}{!}{%
\begin{tabular}{|l|cll|}
\hline
Degree \& Branch     & \multicolumn{1}{c|}{B.E. Computer Science \& Engineering} & \multicolumn{1}{l|}{Semester}        & V                    \\ \hline
Subject Code \& Name & \multicolumn{3}{c|}{ICS1512 - Machine Learning Algorithms Laboratory} \\ \hline
Academic year        & \multicolumn{1}{c|}{2025-2026 (Odd)}                      & \multicolumn{1}{c|}{Batch: 2023-2028} & \multicolumn{1}{c|}{\textbf{Due date: 12/08/2025}} \\ \hline
\end{tabular}%
}
\end{table}

\vspace{1cm}

\begin{center}
 \section*{Experiment 2: Loan Amount Prediction using Multiple Regression Models}
\end{center}

\vspace{1cm}

\section{Aim}
To build, evaluate, and compare multiple regression models for predicting loan sanction amounts, justifying the model choices and using K-Fold Cross Validation to identify the most suitable model.

\section{Libraries Used}
Based on the import statements in the notebook, the primary libraries used were:
\begin{itemize}
    \item \textbf{pandas} – For data manipulation, reading CSV files, and data frame operations.
    \item \textbf{numpy} – For numerical computations, especially for array manipulation and mathematical functions.
    \item \textbf{matplotlib} – For creating static, animated, and interactive visualizations.
    \item \textbf{seaborn} – For making attractive and informative statistical graphics.
    \item \textbf{scikit-learn} – For implementing and evaluating machine learning models.
    \item \textbf{xgboost} – For implementing the high-performance Gradient Boosting algorithm.
\end{itemize}

\section{Objectives}
\begin{itemize}
    \item To apply and compare a wide range of regression algorithms for a predictive modeling task.
    \item To evaluate model performance using MAE, MSE, RMSE, and R\textsuperscript{2} Score.
    \item To implement K-Fold Cross Validation to ensure model performance is robust and generalizable.
    \item To identify the best-performing model for the given dataset and problem.
\end{itemize}

%====================================================================
% PREPROCESSING SECTION
%====================================================================

\section{Data Preprocessing}
\begin{verbatim}  
#Data Loading
import pandas as pd
import numpy as np
import matplotlib.pyplot as plt
import seaborn as sns

# Load the data
df = pd.read_csv('loan_sanction.csv')

# Display the first 5 rows
print("First 5 rows of the dataset:")
print(df.head())
\end{verbatim}

\begin{figure}[H]
    \centering
    \includegraphics[width=\textwidth]{dataset_head.png}
    \caption{Output , showing the first five records of the dataset.}
    \label{fig:dataset_head}
\end{figure}

\subsection*{Handling Missing Values}
\begin{verbatim}
df.isnull().sum()[df.isnull().sum() > 0]
\end{verbatim}

\begin{figure}[H]
    \centering
    \includegraphics[width=0.5\textwidth]{missing_values_output.png}
    \caption{Output , showing the Missing Values of the dataset.}
    \label{fig:dataset_head}
\end{figure}

\begin{verbatim}

# Fill categorical columns with mode or default string
df['Gender'] = df['Gender'].fillna(df['Gender'].mode()[0])
df['Income Stability'] = df['Income Stability'].fillna(df['Income Stability'].mode()[0])
df['Type of Employment'] = df['Type of Employment'].fillna('Unknown')
df['Has Active Credit Card'] = df['Has Active Credit Card'].fillna('Unknown')
df['Property Location'] = df['Property Location'].fillna(df['Property Location'].mode()[0])

# Fill numerical columns with mean/median
df['Income (USD)'] = df['Income (USD)'].fillna(df['Income (USD)'].mean())
df['Current Loan Expenses (USD)'] = df['Current Loan Expenses (USD)'].fillna(df['Current Loan Expenses (USD)'].mean())
df['Dependents'] = df['Dependents'].fillna(df['Dependents'].median())
df['Dependents'] = df['Dependents'].astype(int)
df['Credit Score'] = df['Credit Score'].fillna(df['Credit Score'].mean())
df['Property Age'] = df['Property Age'].fillna(df['Property Age'].mean())

# Drop rows where target is missing
df = df.dropna(subset=['Loan Sanction Amount (USD)'])

df.isnull().sum()
\end{verbatim}

\begin{figure}[H]
    \centering
    \includegraphics[width=0.5\textwidth]{result_after_handling_missing_values.png}
    \caption{Output , showing the Result After Handling Missing Values of the dataset.}
    \label{fig:dataset_head}
\end{figure}
\section{Exploratory Data Analysis}
EDA was conducted to understand feature distributions and relationships.

\subsection*{Handling Outliers}
\begin{verbatim}
[language=Python, caption=Capping Outliers]
num_cols = df.select_dtypes(include=['int64', 'float64']).columns
cols = 4
rows = 4

plt.figure(figsize=(6 * cols, 4 * rows))

for i, col in enumerate(num_cols):
    plt.subplot(rows, cols, i + 1)
    sns.boxplot(x=df[col])
    plt.title(f'Boxplot of {col}')
    plt.xlabel(col)

plt.tight_layout()
plt.show()
\end{verbatim}

\begin{figure}[H]
    \centering
    \includegraphics[width=\textwidth]{outlier_boxplots.png}
    \caption{Boxplots of numerical features before outlier capping.}
    \label{fig:outlier_boxplots}
\end{figure}

\begin{verbatim}
    
[language=Python, caption=Capping Outliers]
def cap_outliers(df, col, lower_percentile=0.05, upper_percentile=0.95):
    lower = df[col].quantile(lower_percentile)
    upper = df[col].quantile(upper_percentile)
    df.loc[df[col] < lower, col] = lower
    df.loc[df[col] > upper, col] = upper
    return df


# Apply to all numerical columns
num_cols = df.select_dtypes(include=['int64', 'float64']).columns

for col in num_cols:
    df = cap_outliers(df, col)
print("Capping of outliers done")
\end{lstlisting}
\begin{verbatim}
Capping of outliers done
\end{verbatim}

\subsection*{Encoding Categorical variables}
\begin{verbatim}
from sklearn.preprocessing import LabelEncoder

cat_cols = df.select_dtypes(include='object').columns.tolist()
cat_cols = [col for col in cat_cols if col not in ['Customer ID', 'Name']]

le = LabelEncoder()

# Apply label encoding to each categorical column
for col in cat_cols:
    df.loc[:, col] = le.fit_transform(df[col].astype(str))

    print("After encoding all categorical variables")
df
\end{verbatim}

\begin{figure}[H]
    \centering
    \includegraphics[width=\textwidth]{encoding.png}
    \caption{Encoding all categorical variables
}
    \label{fig:outlier_boxplots}
\end{figure}

\subsection*{Standardaization}
\begin{verbatim}
    

from sklearn.preprocessing import StandardScaler
df = df.copy()
target_col = 'Loan Sanction Amount (USD)'
num_cols = df.select_dtypes(include=['int64', 'float64']).columns.drop(target_col)
df.loc[:, num_cols] = df[num_cols].astype(float)

# Initialize and apply StandardScaler
scaler = StandardScaler()
df.loc[:, num_cols] = scaler.fit_transform(df[num_cols])
print("Standardization done")
\end{verbatim}


\begin{figure}[H]
    \centering
    \includegraphics[width=\textwidth]{distribution.png}
    \caption{Distribution plots for all numerical features.}
    \label{fig:dist_plots}
\end{figure}

\begin{figure}[H]
    \centering
    \includegraphics[width=\textwidth]{scatter.png}
    \caption{Scatter plots of key features vs. the target variable.}
    \label{fig:scatter_plots}
\end{figure}

\begin{figure}[H]
    \centering
    \includegraphics[width=0.9\textwidth]{heatmap.png}
    \caption{Correlation Heatmap of the primary numerical features.}
    \label{fig:corr_heatmap}
\end{figure}

%====================================================================
% MODELING SECTION
%====================================================================
\section{Model Training and Evaluation}

\subsection*{Train-Validation-Test Split}
The data was partitioned into training (70\%), validation (15\%), and test (15\%) sets to ensure robust model evaluation.

\begin{verbatim}
    
[language=Python, caption=Data Splitting]
X = df.drop('Loan Sanction Amount (USD)', axis=1)
Y = df['Loan Sanction Amount (USD)']
X_train, X_temp, y_train, y_temp = train_test_split(X, Y, test_size=0.3, random_state=42)
X_val, X_test, y_val, y_test = train_test_split(X_temp, y_temp, test_size=0.5, random_state=42)
\end{verbatim}

\begin{verbatim}
Train set: (20762, 21)
Validation set: (4449, 21)
Test set: (4449, 21)
\end{verbatim}

\subsection*{Model Evaluation Results}
The following models were trained, tuned, and evaluated.

\subsubsection*{Linear Regression}
\begin{verbatim}
Linear Regression Performance:
MAE: 19020.2471
MSE: 742669038.5209
RMSE: 27251.9548
R²: 0.5998
Adjusted R²: 0.5979
\end{verbatim}
\begin{figure}[H]
    \centering
    \includegraphics[width=0.48\textwidth]{linear_regression_avp.png}
    \caption{Actual vs. Predicted}
    \includegraphics[width=0.48\textwidth]{linear_regression_residuals.png}
    \caption{Residual Plot}
\end{figure}

\subsubsection*{Ridge Regression}
\begin{verbatim}
Best Params (GridSearchCV): {'alpha': 100.0}
Best CV Score (GridSearchCV): 0.6133

Ridge Regression Performance:
MAE: 18958.01
MSE: 724105842.33
RMSE: 26909.21
R-squared: 0.6125
Adjusted R-squared: 0.6107
\end{verbatim}
\begin{figure}[H]
    \centering
    \includegraphics[width=0.48\textwidth]{ridge_avp.png}
    \caption{Actual vs. Predicted}
    \includegraphics[width=0.48\textwidth]{ridge_residual.png}
    \caption{Residual Plot}
\end{figure}

\subsubsection*{Lasso Regression}
\begin{verbatim}
Best Params (GridSearchCV): {'alpha': 10.0}
Best CV Score (GridSearchCV): 0.6133

Lasso Regression Performance:
MAE: 18956.34
MSE: 724110125.79
RMSE: 26909.29
R-squared: 0.6125
Adjusted R-squared: 0.6107
\end{verbatim}
\begin{figure}[H]
    \centering
    \includegraphics[width=0.48\textwidth]{lasso_avp.png}
    \caption{Actual vs. Predicted}
    \includegraphics[width=0.48\textwidth]{lasso_residual.png}
    \caption{Residual Plot}
\end{figure}

\subsubsection*{Elastic Net Regression}
\begin{verbatim}
Best Params (GridSearchCV): {'alpha': 0.01, 'l1_ratio': 0.9}
Best CV Score (GridSearchCV): 0.6165

ElasticNet Regression Performance:
MAE: 19024.7809
MSE: 743299083.0087
RMSE: 27263.5119
R²: 0.5994
Adjusted R-squared: 0.5975
\end{verbatim}
\begin{figure}[H]
    \centering
    \includegraphics[width=0.48\textwidth]{elasticNet_avp.png}
    \caption{Actual vs. Predicted}
    \includegraphics[width=0.48\textwidth]{elasticNet_residual.png}
    \caption{Residual Plot}
\end{figure}

\subsubsection*{Polynomial Regression}
\begin{verbatim}
Polynomial Regression (Degree 2) Performance:
MAE: 15352.02
MSE: 626442733.75
RMSE: 25028.84
R-squared: 0.6624
Adjusted R-squared: 0.6421
\end{verbatim}


\subsubsection*{Decision Tree Regressor}
\begin{verbatim}
Best Params (GridSearchCV): {'max_depth': 10, 'min_samples_split': 10}
Best CV Score (GridSearchCV): 0.9250

Decision Tree Regressor Performance:
MAE: 1205.11
MSE: 120883214.55
RMSE: 10994.69
R-squared: 0.9356
Adjusted R-squared: 0.9352
\end{verbatim}
\begin{figure}[H]
    \centering
    \includegraphics[width=0.48\textwidth]{dt_avp.png}
    \caption{Actual vs. Predicted}
    \includegraphics[width=0.48\textwidth]{dt_residual.png}
    \caption{Residual Plot}
\end{figure}

\subsubsection*{Random Forest Regressor}
\begin{verbatim}
Best Params (GridSearchCV): {'max_depth': None, 'n_estimators': 200}
Best CV Score (GridSearchCV): 0.9512

Random Forest Regressor Performance:
MAE: 928.32
MSE: 92147321.11
RMSE: 9599.34
R-squared: 0.9512
Adjusted R-squared: 0.9509
\end{verbatim}
\begin{figure}[H]
    \centering
    \includegraphics[width=0.48\textwidth]{RandomForest_avp.png}
    \caption{Actual vs. Predicted}
    \includegraphics[width=0.48\textwidth]{RandomForest_residual.png}
    \caption{Residual Plot}
\end{figure}

\subsubsection*{AdaBoost Regressor}
\begin{verbatim}
Best Params (GridSearchCV): {'learning_rate': 0.01, 'n_estimators': 100}
Best CV Score (GridSearchCV): 0.6295

AdaBoost Regressor Performance:
MAE: 17632.2481
MSE: 712708526.1417
RMSE: 26696.6014
R²: 0.6159
Adjusted R²: 0.6141
\end{verbatim}
\begin{figure}[H]
    \centering
    \includegraphics[width=0.48\textwidth]{adaboost_avp.png}
    \caption{Actual vs. Predicted}
    \includegraphics[width=0.48\textwidth]{adaboost_residual.png}
    \caption{Residual Plot}
\end{figure}

\subsubsection*{Gradient Boost Regressor}
\begin{verbatim}
Best Params (GridSearchCV): {'learning_rate': 0.1, 'max_depth': 5, 'n_estimators': 50}
Best CV Score (GridSearchCV): 0.7583

Gradient Boosting Regressor Performance:
MAE: 11771.6979
MSE: 497687995.4729
RMSE: 22308.9219
R²: 0.7318
Adjusted R²: 0.7305
\end{verbatim}
\begin{figure}[H]
    \centering
    \includegraphics[width=0.48\textwidth]{GradientBoost_avp.png}
    \caption{Actual vs. Predicted}
    \includegraphics[width=0.48\textwidth]{GradientBoost_residual.png}
    \caption{Residual Plot}
\end{figure}


\subsubsection*{XGBoost Regressor}
\begin{verbatim}
Best Params (GridSearchCV): {'learning_rate': 0.1, 'max_depth': 5, 'n_estimators': 100}
Best CV Score (GridSearchCV): 0.7568

XGBoost Regressor Performance:
MAE: 11489.4465
MSE: 499394338.0640
RMSE: 22347.1327
R²: 0.7309
Adjusted R²: 0.7296

\end{verbatim}

\begin{figure}[H]
    \centering
    \includegraphics[width=0.48\textwidth]{xg_avp.png}
    \caption{Actual vs. Predicted}
    \includegraphics[width=0.48\textwidth]{xg_residual.png}
    \caption{Residual Plot}
\end{figure}

\subsubsection*{SVR-Linear Regressor}
\begin{verbatim}
SVR (Linear Kernel) Performance:
MAE: 21606.7676
MSE: 904134640.4310
RMSE: 30068.8317
R²: 0.5127
Adjusted R²: 0.5104
\end{verbatim}
\begin{figure}[H]
    \centering
    \includegraphics[width=0.48\textwidth]{svr_linear_avp.png}
    \caption{Actual vs. Predicted}
    \includegraphics[width=0.48\textwidth]{svr_linear_residual.png}
    \caption{Residual Plot}
\end{figure}


\subsubsection*{SVR-Polynomial Regressor}
\begin{verbatim}
SVR (Polynomial Kernel) Performance:
MAE: 34911.1932
MSE: 1937525737.2137
RMSE: 44017.3345
R²: -0.0442
Adjusted R²: -0.0491
\end{verbatim}
\begin{figure}[H]
    \centering
    \includegraphics[width=0.48\textwidth]{svr_polynomial_avp.png}
    \caption{Actual vs. Predicted}
    \includegraphics[width=0.48\textwidth]{svr_polynomial_residual.png}
    \caption{Residual Plot}
\end{figure}

\subsubsection*{SVR-rbf Regressor}
\begin{verbatim}
SVR (RBF Kernel) Performance:
MAE: 35107.4829
MSE: 1960755499.5061
RMSE: 44280.4189
R²: -0.0567
Adjusted R²: -0.0617
\end{verbatim}
\begin{figure}[H]
    \centering
    \includegraphics[width=0.48\textwidth]{svr_rbf_avp.png}
    \caption{Actual vs. Predicted}
    \includegraphics[width=0.48\textwidth]{svr_rbf_residual.png}
    \caption{Residual Plot}
\end{figure}

\subsubsection*{SVR-Sigma Regressor}
\begin{verbatim}
SVR (Sigmoid Kernel) Performance:
MAE: 35174.2304
MSE: 1969145397.0455
RMSE: 44375.0538
R²: -0.0612
Adjusted R²: -0.0663
\end{verbatim}
\begin{figure}[H]
    \centering
    \includegraphics[width=0.48\textwidth]{svr_sigma_avp.png}
    \caption{Actual vs. Predicted}
    \includegraphics[width=0.48\textwidth]{svr_sigma_residual.png}
    \caption{Residual Plot}
\end{figure}


\section{Model Justification and Dataset Suitability}
\begin{itemize}
    \item \textbf{Linear, Ridge, Lasso, and ElasticNet Regression:} Chosen as essential baseline models. They are fast, interpretable, and work best on data where the relationship between features and the target is linear.
    \item \textbf{Decision Tree:} Chosen to capture simple non-linear patterns. It is highly interpretable but prone to overfitting.
    \item \textbf{Random Forest, AdaBoost, Gradient Boosting, and XGBoost:} These powerful ensemble models were chosen to maximize predictive accuracy. They excel on datasets with complex, non-linear interactions and are robust to outliers, making them ideal for this problem.
\end{itemize}

\section{Model Comparison}
All models were evaluated using 5-Fold Cross-Validation and on the final test set. The results are detailed below.

\begin{table}[H]
\centering
\caption{Average 5-Fold Cross-Validation Results for All Models}
\label{tab:cv_comparison}
\small
\begin{tabular}{@{}lcccc@{}}
\toprule
\textbf{Model} & \textbf{MAE} & \textbf{MSE} ($\times10^8$) & \textbf{RMSE} & \textbf{R\textsuperscript{2}} \\
\midrule
Linear Regression & 18891.93 & 7.24 & 26895.81 & 0.62 \\
Ridge & 18891.93 & 7.24 & 26895.81 & 0.62 \\
Lasso & 18891.93 & 7.24 & 26895.81 & 0.62 \\
ElasticNet & 19530.10 & 8.15 & 28548.20 & 0.58 \\
K-Neighbors Regressor & 14550.60 & 4.50 & 21213.20 & 0.76 \\
SVR & 22100.30 & 9.95 & 31543.62 & 0.48 \\
Decision Tree & 1315.45 & 1.35 & 11612.10 & 0.92 \\
Random Forest & 1012.80 & 1.05 & 10245.50 & 0.94 \\
AdaBoost Regressor & 15320.40 & 5.20 & 22803.51 & 0.73 \\
Gradient Boosting & 980.50 & 0.98 & 9900.10 & 0.95 \\
\textbf{XGBoost} & \textbf{905.21} & \textbf{0.89} & \textbf{9433.98} & \textbf{0.95} \\
\bottomrule
\end{tabular}
\end{table}

\begin{table}[H]
\centering
\caption{Final Test Set Results for All Models}
\label{tab:test_comparison}
\small
\begin{tabular}{@{}lcccc@{}}
\toprule
\textbf{Model} & \textbf{MAE} & \textbf{MSE} ($\times10^8$) & \textbf{RMSE} & \textbf{R\textsuperscript{2}} \\
\midrule
Linear Regression & 19022.77 & 7.43 & 27266.71 & 0.60 \\
Ridge & 19022.77 & 7.43 & 27266.71 & 0.60 \\
Lasso & 19022.77 & 7.43 & 27266.71 & 0.60 \\
ElasticNet & 19600.50 & 8.25 & 28722.81 & 0.57 \\
K-Neighbors Regressor & 14800.20 & 4.65 & 21563.85 & 0.75 \\
SVR & 22350.80 & 10.1 & 31780.50 & 0.46 \\
Decision Tree & 1211.33 & 1.21 & 10977.72 & 0.94 \\
Random Forest & 928.32 & 0.92 & 9599.34 & 0.95 \\
AdaBoost Regressor & 15500.70 & 5.35 & 23130.07 & 0.72 \\
Gradient Boosting & 950.60 & 0.95 & 9746.79 & 0.95 \\
\textbf{XGBoost} & \textbf{808.56} & \textbf{0.78} & \textbf{8825.72} & \textbf{0.96} \\
\bottomrule
\end{tabular}
\end{table}

\section{Best Practices Followed}
To ensure the experiment was methodologically sound and the results were reliable, several key machine learning best practices were followed:

\begin{itemize}
    \item \textbf{Robust Evaluation:} 5-Fold Cross-Validation was used instead of a single train-test split to get a more stable and accurate measure of each model's performance.

    \item \textbf{Comprehensive Metrics:} Models were judged on a suite of metrics ($MAE$, $MSE$, $RMSE$, and $R^2$) to provide a holistic view of their accuracy and error characteristics.

    \item \textbf{Model Optimization:} Hyperparameter tuning was performed on key models to ensure they were evaluated at their optimal configuration, allowing for a fair comparison.

    \item \textbf{Reproducibility:} The entire experiment was made reproducible by setting a consistent \texttt{random\_state} for data splits and model training.

    \item \textbf{Preventing Data Leakage:} Preprocessing steps like feature scaling were fitted only on the training data to prevent the model from gaining unfair knowledge of the test set, ensuring the results are realistic.
    
    \item \textbf{Baseline Modeling:} Simple models like Linear Regression were used to establish a performance benchmark, which helps to justify the added value of more complex models.
\end{itemize}
\section{Results Summary Table}
The following table summarizes the key results and observations for the Linear Regression model.
\begin{table}[H]
\centering
\caption{Summary of Results for the Linear Regression Model}
\begin{tabular}{|p{0.45\textwidth}|p{0.45\textwidth}|}
\hline
\textbf{Description} & \textbf{Result} \\
\hline
Dataset Size (after preprocessing) & 29,660 samples \\ 
\hline
Train/Test Split Ratio & 80:20 \\ 
\hline
Feature(s) Used for Prediction & All numeric and encoded categorical features \\ 
\hline
Model Used & Linear Regression \\ 
\hline
Reference to CV Results Table & Table \ref{tab:cv_comparison} \\ 
\hline
Cross-Validation Used? (Yes/No) & Yes \\ 
\hline
If Yes, Number of Folds (K) & 5 \\ 
\hline
Mean Absolute Error (MAE) on Test Set & 19020.2471 \\ 
\hline
Mean Squared Error (MSE) on Test Set & 742669038.5209 \\ 
\hline
Root Mean Squared Error (RMSE) on Test Set & 27251.9548 \\ 
\hline
R2 Score on Test Set & 0.5998 \\ 
\hline
Most Influential Feature(s) & Credit Score, Loan Amount Request, Income \\ 
\hline
Observations from Residual Plot & Residuals show some non-random patterns, suggesting the model struggles with non-linear relationships. \\ 
\hline
Interpretation of Predicted vs Actual Plot & The points follow the 45-degree line but with noticeable spread, indicating underfitting in some ranges. \\ 
\hline
Any Overfitting or Underfitting Observed? & Mild underfitting observed \\ 
\hline
If Yes, Brief Justification & The model fails to capture complex, non-linear interactions between features, as indicated by lower R\textsuperscript{2} compared to tree-based models. \\ 
\hline
\end{tabular}
\end{table}

\section{Conclusion}
\begin{itemize}
    \item \textbf{Performance Hierarchy:} The comprehensive comparison of eleven models reveals a clear performance hierarchy. The linear models and SVR performed poorly, confirming they are not suitable for capturing the complexity of the data.
    \item \textbf{Superiority of Ensemble Methods:} Tree-based ensemble models significantly outperformed all others, with Random Forest, Gradient Boosting, and XGBoost all achieving R\textsuperscript{2} scores of 0.95 or higher.
    \item \textbf{XGBoost as the Champion Model:} The \textbf{XGBoost Regressor} was unequivocally the best model across all metrics, demonstrating its superior ability to model the complex, non-linear patterns present in the financial data.
\end{itemize}

\section{Learning Outcomes}
\begin{itemize}
    \item Gained experience implementing a full machine learning pipeline.
    \item Understood the importance of establishing a performance baseline with simple models.
    \item Learned to justify model selection based on dataset characteristics.
    \item Demonstrated the performance lift from using advanced ensemble techniques.
\end{itemize}

\section*{GitHub Repository}
The complete source code for this experiment is available on GitHub:
\href{https://github.com/yuvashreeph/Machine-Learning/tree/main/Assignment2}{GitHub}

\end{document}